\section{Support for Community Science}
\label{sec:community_science}

Rubin Observatory's model to support its large, diverse science community to access and analyze \gls{DP2} focuses on providing documentation, tutorials, and live virtual learning opportunities alongside an open platform for issue reporting and resolution that enables crowd-sourced solutions.
These support resources are all available via the main hub of Rubin's "For Scientists" website\footnote{\nolinkurl{https://rubinobservatory.org/for-scientists}}.

The Data Release documentation for \gls{DP2}\footnote{\nolinkurl{https://dp2.lsst.io}} provides an overview of the \gls{LSSTCam} observations, data quality summary plots, and descriptions of the contents and access methods for all catalog and image datasets.
All tutorials \citep[][]{10.11578/rubin/dc.20250909.20} that demonstrate \gls{DP2} access via the Rubin Science Platform are delivered in the Data Release documentation, and executable versions of the Jupyter notebook tutorials are also available from a drop-down menu within the Notebook Aspect.

Early \gls{DP2} was released on the first day of the annual Rubin Community Workshop, July 27 2026.
A dedicated "Get started with DP2" session was added to the Wednesday agenda.
Four DP2-onboarding virtual sessions were run in the first week of August, a mix of presentations, tutorials, and time for issue resolution.

The Rubin Community Forum\footnote{\nolinkurl{https://community.lsst.org/}} is the open platform for user issue reporting and resolution.
Rubin staff monitor all new topics posted in the Support category and ensure that all requests for help get a response and a solution.
The open nature of the platform also enables peer-to-peer support that avoids knowledge bottlenecks and hastens the time to resolution.
